\section{结论}\label{sec:conclusion}
本文研究路网拓扑已知、部分节点不参与动态信号训练时，窗口内 OS 变化下的归纳时空预测，并提出 ALFNet。ALFNet 以逐步的观测角色、状态进度和观测掩码描述 OS：SAHC 为两个分支提供同一份历史，ALFTemporal 从跨节点历史生成条件特征并经 FiLM 调制图编码状态，ALFGraph 依据生命周期上下文组合图传播算子并解码，输出端再施加由末步观测角色控制的有界修正。本文的贡献不在于提出新的基础算子，而在于以 OS 组织补全、传播与上下文的计算方式。

实验证据与引言的论点逐一对应。总体与场景评价表明，与八个比较方法（共同协议下的适配版本）相比，ALFNet 在四个数据集的 LifeAvg 三项指标上均为最优，Join 与 Fault 场景的优势在四个数据集上一致；组件消融表明，跨节点时间上下文与时间条件分支的作用最大，OS 条件化传播与状态进度带来较小但方向基本一致的改进；在预测目标与非干预输入保持不变的配对干预中，加入或删除本地历史、改变可见读数的时间位置都伴随相应的预测响应，但布局上的优势在验证集与测试集之间并不一致。这些结果与本文的核心主张一致：OS 应同时作为输入有效性的描述和预测计算的条件，而不能仅由训练身份或掩码比例刻画；不过，OS 条件化本身的独立作用较小，仍需整体移除 OS 条件的对照与多次重复训练加以确认。

结果的边界同样明确。未取得第一的结果集中在本地历史接近缺失的端点：两个流量数据集的 Unseen 场景、PEMS-BAY 的 Cold 与 Outage 阶段，以及 PEMS04 的完全中断干预；PEMS08 的部分中断与仅保留早期读数的 Early 布局也劣于代表基线。ALFNet 的优势在本地证据部分存在时最为稳定，Unseen 场景下的泛化仍需改进。此外，实验使用单一随机种子，部分历史配置的保留曾参考测试结果，因此测试结果是既定研究流程下的最终评估，而非全程未触碰的独立确认；比较方法是共同输入与评价边界下的适配版本，SLPF 与 STBP 未执行其原始的多阶段或持续学习流程，基线未像 ALFNet 一样经过针对性调参，比较中也未包含 Graph WaveNet、AGCRN 等自适应图方法以及历史均值等简单基线；速度与流量数据在窗口插值、解码残差倍率和中心化上使用不同的固定设置；机制分析在已知节点的验证窗口上进行，且只有点估计；评价中的观测角色与状态进度由协议精确给出，实际应用中由接入记录与故障检测带来的延迟或误判没有评估。本文只在交通速度与流量数据上验证，结论能否推广到其他传感网络，取决于能否获得可靠的接入与故障状态。后续工作可着重改进本地历史完全缺失时对空间证据的利用；当前解码阶段固定使用末步生命周期上下文，未来也可引入对预测窗口内观测角色演变的建模。
